\subsection*{Continuous visual inputs carry inherited contextual knowledge}
We removed the vocabulary lookup from a pretrained Chinese extractive reader and supplied continuous vectors computed from a pictured passage and question. A first regression from ink images recovered 48/64 answers on the previously specified new scenes, close to the native text reader's 49/64. It recovered six complete four-question scenes, also matching the native reader. Changing only the typeface to a held-out bold face reduced accuracy to 1/64. The mean cosine similarity of 0.9991 between the fitted image representations and their targets therefore gave a misleading impression of robustness.

A visual metric trained on 54,956 glyph renderings brought bold accuracy to 41/64 with three complete scenes, while regular print gave 46/64 with six complete scenes (Fig.~\ref{fig:font-invariant-reader}a). The same weights gave 1/64 in the unseen Kai family. All 64 questions and names were renderable in every face. Exchanging the two questions about a regular-print passage reduced the score to 14/64 and zero complete scenes. On Chinese XWinograd, the new metric gave 79/504 in regular print, 81/504 in bold and 3/504 in Kai. The native text reader gave 73/504. Both registered cross-typeface criteria failed. These experiments measure transfer of language knowledge from a pretrained encoder into a visual interface. They do not measure the cost of acquiring that knowledge from pictures.

The existing image writer's learned decision could invoke this reader or compose an answer recurrently. The combined responder gave 263/387 natural person answers and 69/131 correct pairs, up from 148 and 19. On complete DRCD passages it gave 2,698/3,505 exact answers, compared with 1,344. Prior exposure of the inherited encoder to benchmark material was not audited. On the 16 investigator-written new scenes in regular print it gave 46/64 and six complete scenes, compared with 55/64 and ten for local Qwen3-8B receiving text. With historical names it composed 30/32 familiar descriptions. It copied the correct historical participant on 39/96 four-name inventory questions, with no inventory solved across every rotation. Familiar reordered statements remained at 344/404, unfamiliar wording at 8/404, and open DuReader answers at 0/571. The complete numeric weights and writing memory occupy 395.1 MB in storage, excluding software. These reading results leave general independent writing and matched-quality compute savings open.
